\begin{afterword}
\summaryhead

The post is a fable about a species whose deepest passion is sorting pebbles into correct heaps. The Pebblesorters cannot say why a heap is correct, but their history looks like progress. A great leader's heap of 91 pebbles is later judged incorrect and scattered, and a war over heaps of 1957 ends when a philosopher shows a heap of 103 beside a heap of 19. The reader can work out what the characters cannot: the correct heaps are the prime ones.

Some philosophers, the Heap Relativists, hold that there is no real progress and nothing that makes a heap correct. They warn that an artificial intelligence might favour other heaps, and propose strapping bombs to computers. Most Pebblesorters find the warning absurd: a superintelligence would see at once which heaps are correct, and would rewrite a utility function that told it otherwise. After all, smarter minds have always made smarter heaps. The story ends by asking why that trend would break.

\responsehead

As fiction the parable works because it withholds its rule. The reader stands where an outsider would, able to see a fixed criterion that the characters use and cannot state, and the clues (``But\ldots\ 13!'', a war ended by the factors of 1957) are fair and funny. The satire falls on both camps: the relativists are mistaken about progress and absurd in their remedy, and the majority are mistaken about AI.

As argument, the story makes two points, and both hold within it. The first answers the Heap Relativists. Disagreement over hard cases and change over time are compatible with a fixed criterion: the reader can see that deciding whether a large heap is prime is simply hard. This is the standard objectivist reply to the argument from disagreement, and the complaint that relativism cannot account for progress is also standard. But the progress the story shows is progress in working out a criterion the species already holds, which is the kind a relativist can accept, as the Stanford Encyclopedia notes. The story answers the Heap Relativists' denial of progress, not relativism in general, and the post claims no more.

The second point is the one the story is built for. The majority infer from a trend among Pebblesorters, smarter minds make smarter heaps, that any intelligence will favour the same heaps. The story shows why the inference fails: the trend holds because every mind in it shares one criterion. An AI could tell at a glance which heaps are prime without caring. Omohundro had argued earlier the same year that such an agent would protect its utility function rather than rewrite it.

The story illustrates this point rather than arguing it. That an AI need not share the criterion is part of the setup, since the reader knows that primality has no special hold on a mind. The argument for the general claim is in ``No Universally Compelling Arguments,'' earlier in this book. Likewise, the post does not say that human values resemble the Pebblesorters' criterion; it leaves that application to the reader and to later posts.

In short: a well-made fable that shows two things clearly, that disagreement and change are compatible with a fixed criterion, and that a criterion a species shares need not be shared by other minds; the second is stipulated by the story and argued elsewhere, and the relativists it answers are those who deny progress.
\end{afterword}
